\section{问题分析}

\iffalse
\cite{weiss_2005_dpd,penna_2013_doa,bayram_2016_bearing,russell_2020_doa,zhao_2013_placement,moreno_2013_placement,galceran_2013_cpp,cabreira_2019_uavcpp,vansteenwegen_2011_op,gunawan_2016_op,li_2020_crlb,li_2002_tracking}
\fi
四问共享同一套观测、动作与时钟，差别只在信息集与完备性证书。
若把问题3、问题4理解成尽快找到大约十个源，会在尚未证明频道判空时提前退出。
若把问题1的直径当成最小覆盖圆直径，后续用半径 $20\,\mathrm{m}$ 的光学清除会留下死角。
本节先把这套共享结构写清楚，再按问拆开任务。
流程图按问放入对应求解章，避免在分析里把四张流程叠在一起。

测向接口返回三值。记检测点为 $q$、源位为 $g$、距离 $d=\|q-g\|$。当 $q$ 落在该源有效覆盖内且 $5<d\le r_i$ 时，返回 $\mathtt{direction}$ 及示向度 $\hat\theta$；当覆盖内且 $d\le 5$ 时返回 $\mathtt{near}$，不得把缺失的角度填成 $0^\circ$；其余情形返回 $\mathtt{no\_signal}$。问题3在已经确认某频道仍有未清除全向源时，可以把 $\mathtt{no\_signal}$ 读成 $d>r_i\ge 1000$；一旦存在性未知，就必须保留“该频道根本无源”的分支。问题4还要加上第三分支：定向扇外。射频阴影里仍可能 $d\le 20$，因而光学清除与射频覆盖必须解耦。同点同频道重复测量仍支付 $5\,\mathrm{s}$ 检测时间，但误差样本并不刷新。

清除成功当且仅当指定频道存在未清除源且欧氏距离不超过 $20\,\mathrm{m}$。同一源只能成功一次；二次清除不计新的 $K$，却仍支付光学时间。上一合法位置为 $p$、测向机频道为 $c$，本次目标为 $q$、指定频道为 $k$ 时，虚拟耗时为
\begin{equation}
\Delta T_{\mathrm{measure}}=\frac{\|q-p\|}{5}+\mathbf{1}_{k\neq c}+5,\qquad
\Delta T_{\mathrm{clear}}=\frac{\|q-p\|}{5}+3+2\mathbf{1}_{\mathrm{success}}.
\end{equation}
$\mathtt{/clear}$ 不改写测向机频道 $c$，只有 $\mathtt{/measure}$ 会切换频道。$\mathtt{/enter}$ 与 $\mathtt{/exit}$ 不推进虚拟时间。被拒绝的指令不得覆盖时钟。任务效率使用虚拟时间 $T$，程序运行时间是 $\mathtt{/enter}$ 到结束的墙钟（上限 $1200\,\mathrm{s}$），测试窗口另有 $1500\,\mathrm{s}$；三套时钟不得混用。虚拟时间上限 $360000\,\mathrm{s}$ 只约束模拟器账本，不是现场必须用满的预算。

测向定位本身是经典测向交会问题。Weiss 与 Amar 把多站测向直接写入位置域似然\upcite{weiss_2005_dpd}；认知无线电中仅用到达方向也可完成主用户定位\upcite{penna_2013_doa}。本题接口不提供场强或距离，不能做距离回归，只能使用示向度与覆盖几何。Bayram 等人讨论如何主动采集方位数据\upcite{bayram_2016_bearing}，Zhao 等人则从几何上分析传感器相对目标的最优构型\upcite{zhao_2013_placement}。无卫星条件下的无人机也可仅用到达方向完成协同定位\upcite{russell_2020_doa}，但本题只有单机、单频道，不能靠多站瞬时交会。三维到达角的误差下界有闭合形式讨论\upcite{li_2020_crlb}，本题忽略高度差，对象退化为平面闭楔形。这些工作给出的是统计最优或渐近方差，本题误差被硬性限制在 $\pm 1^\circ$，同一点重复测量不是新样本，因而更自然的对象是闭楔形交集，而不是无约束的最小二乘交点。

问题1的定位区域 $R$ 定义为各检测点前向闭楔形（示向度 $\pm 1^\circ$）的交集。
两点时它是题面图2中的四边形，多点时是多个楔形的公共部分。
直径定义为
\begin{equation}
\mathrm{diam}(R)=\sup_{x,y\in R}\|x-y\|,
\end{equation}
对凸多边形它等于顶点最远对。
以该直径为直径的圆一般不能覆盖 $R$：等边三角形的直径是边长，而以边为直径的圆过另外两个顶点却把第三个顶点留在圆外。
Jung 定理给出的最小覆盖圆半径约为边长的 $1/\sqrt{3}$，严格大于直径之半。
平面凸集上还有不等式 $\mathrm{diam}(R)/2\le\rho_\ast\le\mathrm{diam}(R)/\sqrt{3}$：左端是任意覆盖圆的平凡下界，右端由 Jung 定理给出。
等边演示取到右端量级，两点楔形演示则取到左端附近。
若后续用 $R$ 选清除点，必须检查该点到 $R$ 的豪斯多夫半径不超过 $20\,\mathrm{m}$，不能用 $\mathrm{diam}(R)/2$ 代替，也不能把 $\mathrm{diam}(R)\le 40$ 写成可清。
题面不保证 $R$ 有界非空，空集、无界楔与退化线段都必须作为合法输出，而不能把无界写成一个有限数字。

问题2只有一个全向源的一次示向度，真位置与 $r_i$ 均未知。
第二点必须落在“可能收到信号”的几何范围内，又不能用未知 $g$ 来排点。
传感器布置文献强调基线长度与夹角对定位误差的影响\upcite{zhao_2013_placement,moreno_2013_placement}，但那些指标通常依赖未知真值。
本文改用外包可行多边形上的最坏后验半径作为排序指标 $J(q)$，并明确它只是候选集推荐，不是清除证书。
理论扇半径 $\rho_{\mathrm{th}}=r_{\max}/(2\cos\varepsilon)\approx 750.1142460329307\,\mathrm{m}$ 只用来标定核心尺度。
演示外包半径约 $750.2285094687356\,\mathrm{m}$，核心可观测半径约 $249.77149053126436\,\mathrm{m}$。
演示推荐点约 $(719.0070731523278,634.1607723730604)$，$J^\ast\approx 43.0312146778083\,\mathrm{m}$，已经大于光学半径 $20\,\mathrm{m}$，说明第二点之后通常还要累计测向或光学网格。
开发集二十例上 $J^\ast$ 中位数 $42.75989\,\mathrm{m}$、最大值 $46.20539\,\mathrm{m}$，全部高于 $20\,\mathrm{m}$，与演示同一结论。

问题3、问题4是带覆盖约束的在线路径规划。
传感器网络文献把检测、分类与跟踪分成层次\upcite{li_2002_tracking}；本题只有单机、单频道，层次被压成“发现覆盖 $\to$ 局部交会 $\to$ 光学清除 $\to$ 频道判空”。
覆盖路径规划的综述把“保证扫过工作区域”与“尽量缩短路径”分开处理\upcite{galceran_2013_cpp,cabreira_2019_uavcpp}；定向越野问题则在时间预算下选择有收益的顶点\upcite{vansteenwegen_2011_op,gunawan_2016_op}。
本题的字典序目标是先 $K=N$、再最小化 $T$，因此完备性证书优先于启发式加速。
收益顶点在本题里是未知频道上的首次 $\mathtt{direction}$ 或 $\mathtt{near}$，漏清不可用超时换取。
“已找到 $10$ 个”“某点二十个频道都无信号”都不是题设保证的全清判据。
问题3用原点加半径 $1200\,\mathrm{m}$ 的正六边形七点作为发现覆盖，同站批量扫描未决频道。
问题4在未知朝向下必须改用能同时覆盖 $180^\circ$ 半平面与最小接收半径的有限点集：SQUARE81 共 $81$ 点、步长 $600\,\mathrm{m}$，HEX37 共 $37$ 点、步长 $800\,\mathrm{m}$，并在全部点上对未决频道测得 $\mathtt{no\_signal}$ 才可判空。
光学清除另用边长 $28\,\mathrm{m}$ 的方格覆盖连续后验，半对角线 $14\sqrt{2}<20\,\mathrm{m}$，与发现格点不是同一证书。

信息结构上，正式测试永不揭示 $N$，论文不得把 $K$ 写成百分之百清除率。演练测试结束时返回 $N$，可用于计算 $K/N$，但场景随机、局间不配对，不能当作正式成绩，也不能单独支撑“策略改进百分之多少”的因果断言。离线几何模型可以固定种子做配对，却不经过 HTTP 接口，也不经过模拟器的墙钟与登录校验。后文凡引用演练或离线几何模型的数字，均在表注或正文中标明口径。写作时问题3、问题4的正式测试均未执行，表1保持空缺。

按问拆开后，求解顺序是固定的。
首先用闭楔形交给出定位区域及其直径、Jung 半径，并回答直径圆一般不能覆盖。
再在单次示向度的凸外包上构造分层候选，用 $J(q)$ 排序第二点，同时声明该指标不是清除证书。
接着对全向源给出七点发现巡回与 GREEDY\_FAST 在线策略，用 SAFE 走廊或边长 $28\,\mathrm{m}$ 的光学方格覆盖连续后验。
最后对混合源证明 SQUARE81 与 HEX37 的覆盖命题，累计全部正向观测，并在冻结 Hamilton 骨架上做有界插入。
四问的流程图分别放在第5至第7节，分析章只保留这条逻辑骨架。
参数全部冻结在 $28$ 项表中，求解器不另写一套常数；数值包络 $1.005^\circ$ 只吸收两位小数量化，不写入参数表。
问题3非正式十局均值 $341\,\mathrm{s}$、问题4离线配对 $836.91\,\mathrm{s}$ 与 PREFIX\_A 演练 $822.38\,\mathrm{s}$ 分属三条口径，分析章只标明口径，不把它们平均成现场成绩。

参数全部冻结在 $28$ 项表中，求解器不另写一套常数。圆域半径 $1800\,\mathrm{m}$、源数上下界 $10$ 与 $16$、频道 $1$ 至 $20$、接收半径 $[1000,1500]$、误差 $1^\circ$、近区 $5\,\mathrm{m}$、清除 $20\,\mathrm{m}$、速度 $5\,\mathrm{m/s}$、检测 $5\,\mathrm{s}$、切换 $1\,\mathrm{s}$、光学 $3\,\mathrm{s}$、激光 $2\,\mathrm{s}$、墙钟上限 $1200\,\mathrm{s}$、测试窗口 $1500\,\mathrm{s}$、虚拟上限 $360000\,\mathrm{s}$、正式额度每问三次，均直接来自题面与附件，不在正文中改数。数值包络 $1.005^\circ$ 只作为量化约定出现在后问实现里。后文引用的演练均值与离线配对均值，来源分别是问题3主线说明、配对汇总与演练记录，口径在首次出现处标明。

跨问联动可以写成四条，避免后文各写一套定义。
第一，问题1的闭楔形与顶点最远对是后问累计定位的原子操作：每来一次 $\mathtt{direction}$，就裁进两条半平面，并更新直径与包围半径。
第二，问题2的 $J(q)$ 只在单源、单次示向度下排序第二点；问题3发现源后调用同一几何，但巡回途中改用固定的沿向 $500\,\mathrm{m}$、横向 $150\,\mathrm{m}$ 偏移，以免为理论最优点破坏七点顺序。
第三，问题3的判空在已确认该频道仍有全向源时可以把无信号读成距离圆盘，问题4禁止这条推理。
第四，光学方格边长 $28\,\mathrm{m}$ 来自半对角线小于 $20\,\mathrm{m}$，与问题1不能用直径之半代替包围半径是同一几何：覆盖连续区域必须用能盖住格子的圆，而不是用集合宽度除以二。

信息集上还有两处容易写错。
演练结束弹窗里的 $N$ 是事后核对用的，在线策略若把它写回状态机，就等于在正式测试里偷看答案。
离线几何模型可以固定种子、固定端点误差，用来做配对与回归，但它不经过登录、倒计时与墙钟校验，不能替代正式三次。
后文凡出现 $30/30$、$10/10$ 或具体 $T/K$，都会同时写口径：开发集、离线配对、非正式演练或正式缺失。
把不同口径的数字平均成一个现场成绩，是比算错直径更严重的交付错误。
问题3非正式十局均值 $341\,\mathrm{s}$、问题4离线配对 $836.91\,\mathrm{s}$ 与 PREFIX\_A 演练 $822.38\,\mathrm{s}$ 分属三条口径，禁止平均。
